%	RESUMEN
\chapter{Resumen}

Estas instrucciones servirán de guía para la preparación de los trabajos que se presentarán como requisito en el proceso de graduación de la Materia Integradora de la Unidad Académica. El resumen deberá contener entre 150 a 200 palabras, incluyendo los siguientes cuatro componentes, descritos de forma concisa, comprensible y redactado en estilo impersonal (tercera persona): 1) se empieza con una breve introducción, objetivos, hipótesis y justificación del proyecto descrito en tiempo presente; 2) desarrollo del proyecto, donde se describirán brevemente los materiales, equipos, técnicas, normas etc. utilizadas en el proyecto. Esta sección se redacta en tiempo pasado; 3) continúe con los resultados donde se describen de forma concisa solo los principales y más novedosos resultados escritos en tiempo pasado; 4) finalmente, se presentan las conclusiones generales del proyecto en tiempo presente. Además, deberá incluir al menos 4 palabras clave al final del documento. Todo el resumen se presentará en un sólo cuerpo. Utilice el contador de palabras del procesador de texto para asegurarse del tamaño del documento. 

\textbf{Palabras Clave:} Formato, Proyecto Integrador, etc. (mínimo 4 y máximo 5 palabras).
Procure que las palabras claves que aquí coloque no repitan las que ya se encuentran incluidas en su título.


%	ABSTRACT
\chapter{Abstract}
These instructions will serve as a guide for preparing the papers that will be submitted as a requirement for graduation in the Integrative Subject of the Academic Unit. The abstract should contain between 150 and 200 words, including the following four components, described concisely, comprehensibly, and written in an impersonal style (third person): 1) Begin with a brief introduction, objectives, hypothesis, and justification of the project, described in the present tense; 2) Project development, where the materials, equipment, techniques, standards, etc., of the project will be briefly described. This section is written in the past tense; 3) Continue with the results, where only the main and most novel results are concisely described, written in the past tense; 4) Finally, present the general conclusions of the project in the present tense. In addition, include at least four keywords at the end of the document. The entire abstract should be presented in a single document. Use the word counter in your word processor to ensure the document size is correct.

\textbf{Keywords:} Gesture Detection, Automatic Feedback, Video Processing, Artificial Intelligence, Data Analysis.

\newpage